% Notes of the Montreal seminar, 1970 (cote 129, typed part): archive cote 129, site page 57 (PDF page 58).
% typescript. Draft compared: none.
% AI-transcribed by Claude Opus 5.5 (claude-opus-5-5, adaptive thinking, effort high), Message Batches API, batch mtl129_full (msgbatch_01AGnDKHwy6ATCM7ehUa6263), 2026-11-13; not yet checked by eye.
\scanpage{129/57}
\mspage{\hand{III-2}}

2. \uline{Platitude et critère de représentabilité}

2.1. Soit $G$ un $\Lambda_n$-module. On dit que $G$ est \uline{plat sur} $\Lambda_n$ si
le produit tensoriel par $G$ au dessus du faisceau constant $\Lambda_n$ préserve
les monomorphismes; on a alors le critère suivant:

2.2. \uline{Proposition}: Soit $G(n)$ un groupe annulé par $p^n$. On a les
équivalences:

(a) $\forall i \quad 1 \leq i \leq n, \quad G(i)$ est un $\Lambda_i$-module plat

(b) $G(n)$ est un $\Lambda_n$-module plat

(c) $\forall i, \quad 1 \leq i \leq n-1, \quad p^{n-i}\, G(n) = G(i)$

(i.e. l'application $G(n) \xrightarrow{p^{n-i}\mathrm{Id}} G(i)$ est un épimorphisme)

(d) $\exists i, \quad 1 \leq i \leq n-1,$ tel que $p^{n-i}G(n) = G(i)$

(e) $p\, G(n) = G(n-1)$

(f) Le morphisme canonique
\[ \gamma\colon \Lambda_1[T]/T^n \underset{\Lambda_1}{\otimes} G(n)/p\, G(n) \to gr_p G(n) = \bigoplus_{i=0}^{i=n-1} p^i G(n)/p^{i+1}G(n) \]
(où $\gamma^i\colon G(n)/p\, G(n) \to p^i G(n)/p^{i+1}G(n)$ est induit par la
multiplication par $p^i$) est un isomorphisme.

(g) Le morphisme
\[ \beta\colon gr_p G(n) \to \Lambda_1[T]/T^n \underset{\Lambda_1}{\otimes} G(1) \]
(où $\beta^i\colon p^i G(n)/p^{i+1}G(n) \to G(1)$ est induit par la multipli-%
cation par $p^{n-i-1}$) est un isomorphisme.

\uline{Démonstration}:

(b) $<=>$ (c) Pour tout $i$, $1 \leq i \leq n-1$, on a des triangles commutatifs

\mspage{\hand{57}}
